\section{Uniform geometric pressure by finitely many transverse tests}
\label{sec:v100-pressure}
The fixed-table conclusion in Corollary~\ref{cor:v99-fixed-radius}
is not interchanged with a supremum over tables. We prove a uniform
bound directly. The argument uses expansion and distortion, but no
mixing time, invariant Cantor rectangle or parameter-dependent
spectral gap.

We use standard collision coordinates $(s,\varphi)$ in this section;
$s$ is physical arclength and $p=\sin\varphi$. Jacobians along curves
refer to Euclidean arclength in these coordinates. The index of a
homogeneity strip is denoted by $k$; its width is $O(|k|^{-3})$.
Let $\mathcal M_0^{m,H}(R)$ be the homogeneous continuity partition
of \cite[(2.4)]{BDPressure}, and set
\begin{equation}\label{eq:v100-Q}
 Q_m(R)=\sum_{A\in\mathcal M_0^{m,H}(R)}
               \sup_{x\in A\cap M_R'}J^sT_R^m(x),\qquad Q_0(R)=1.
\end{equation}
These are geometric partition sums, without the section, guard or
four-record twist. They are used only as positive upper bounds for
the original exact-label fibers through Theorem~\ref{thm:v99-pressure-comparison}.

\begin{lemma}[Uniform geometric constants]\label{lem:v100-geometric}
The given circular family admits one choice of cone bounds,
homogeneity threshold, short-curve length and constants with the
following properties.

For every homogeneous short cone-stable curve $W$, let
$\mathcal G_j(W)$ be its homogeneous inverse-image components,
subdivided at the chosen length scale. Then
\begin{equation}\label{eq:v100-curve-growth}
 \sum_{V\in\mathcal G_j(W)}\sup_V J_VT_R^j\le C_g,
                                      \qquad j\ge0.
\end{equation}
On one homogeneous backward continuity cell, the ratio of the
inverse Jacobians on any two short cone-stable curves is between
$C_d^{-1}$ and $C_d$, independently of the iterate.

For $A\in\mathcal M_0^{m,H}(R)$, let $B_{m-1}(A)$ be the cell of
$\mathcal M_{-m+1}^{0,H}(R)\vee H$ containing $T_R^{m-1}A$.
There are $\delta_*>0$ and $c_{\rm long}>0$, independent of $R,m$, such that
\begin{equation}\label{eq:v100-long-cell}
 c_{\rm long} Q_m(R)\le
 \sum_{\substack{A\in\mathcal M_0^{m,H}(R)\\
             \operatorname{diam}^u B_{m-1}(A)\ge\delta_*}}
                          \sup_A J^sT_R^m.
\end{equation}
The geometric constant $c_{\rm long}$ is unrelated to the unchanged
section normalization $c_*$.
\end{lemma}
\begin{proof}
The common geometric bounds are a positive minimum free flight,
a finite maximum free flight, positive lower and upper curvatures,
and bounded third derivatives of the circular boundaries. The
number of possible one-flight lattice displacements is uniformly
finite. Physical arclengths can be identified with fixed collision
circles by uniformly bounded linear changes of coordinate. Thus
\cite[Theorem 2.5 and Remark 2.7]{DZUniform} applies with one set of
geometric constants. Its uniform growth estimate, Lemma 3.2(b), is
\eqref{eq:v100-curve-growth}. The whole-cell distortion statement is
\cite[Lemma 2.2]{BDPressure}; its proof sums the common hyperbolic
contraction and homogeneous distortion bounds and therefore gives
one $C_d$ on this family. It compares different stable curves, not
merely two points on the same curve.

For clarity, the uniformity needed in the long-cell lemma is a
finite geometric choice, rather than a compactness assertion about
$\sup_m Q_m(R)$. Use \cite[Lemmas 3.3, 3.5 and 3.6]{BDPressure}
only at potential parameter one and additional potential zero.
In their proof choose the one-step small-curve contraction
$\vartheta<1$ and its prefactor $C_0$ from the common growth bounds.
Choose an integer $n_1$ so that the resulting geometric loss series
$C_0\sum_{j\ge1}\vartheta^{jn_1}$ is smaller than the fixed loss
fraction used in Lemma 3.5. This choice is independent of $R$.
The short-curve threshold for the first $2n_1$ iterates is then chosen
using the same one-step expansion and homogeneous tail. Only finitely
many iterates have been specified at this stage. Their branchwise
H\"older moduli and complexity bounds are uniform: the free-flight
formula has the same square-root contact singularity, bounded
geometric coefficients and a finite common displacement alphabet.
The identical construction applies to $T_R^{-1}$.

In the proof of Lemma 3.6 choose the smaller of the forward and
backward length thresholds, and then choose $n_2$ with
$C\Lambda^{-n_2}$ below that threshold, where $\Lambda>1$ is the
common cone expansion. Partition a backward cell by its most recent
ancestor having the prescribed unstable length. The branches which
have not yet reached that length are charged to the preceding
small-curve geometric series. Applying distortion and summing that
series leaves a fixed positive fraction of the full pressure on
cells with a long unstable curve. This is precisely the inequality
of that lemma, with its $\delta_2/3$ renamed $\delta_*$. Its remaining
finitely many initial iterates are handled by reducing the common
length threshold: the homogeneous tails have uniformly summable
weights, and the remaining finite physical branches have the same
uniform geometric bounds. This gives \eqref{eq:v100-long-cell}.
No step of this choice uses Proposition 3.14 of that paper, its
mixing argument, or its table-dependent return time. The section normalization is not used in this geometric lemma.
\end{proof}

\begin{lemma}[Finite testing of long backward cells]
\label{lem:v100-tests}
There are integers $N_*,D_*<\infty$ and, for each $R$, homogeneous
short cone-stable curves $W_1(R),\ldots,W_{N_*}(R)$ such that every
cell $B$ in \eqref{eq:v100-long-cell} meets at least one test curve.
The same $N_*$ works for all $R$ and all $m$. At most $D_*$ different
forward cells $A$ have a given $B_{m-1}(A)=B$.

Distinct backward cells assigned to one test can be represented by
distinct members of $\mathcal G_{m-1}(W_j)$. The representatives are
chosen for summation only and introduce no new physical stopping rule.
\end{lemma}
\begin{proof}
A homogeneous cone-unstable arc of length $\delta_*$ has a coordinate
projection of length at least $c\delta_*$. Since its slope is bounded
above and away from zero, it cannot lie in a homogeneity strip of
width smaller than $c'\delta_*$. Hence only finitely many strip indices
can occur. Work in finitely many overlapping lifted collision charts.
Choose a slope strictly inside the common stable cone. The affine
coordinate transverse to lines of that slope increases along every
cone-unstable arc at speed bounded below by a positive constant.
A grid of transverse levels with mesh less than $c''\delta_*$ must
therefore meet every such arc in its interior. Split these grid lines
into overlapping short pieces and intersect with the finitely many
relevant homogeneity strips. This produces at most $N_*$ tests.
The overlap removes chart-end ambiguities. A generic translation of
the grids avoids coincidence with true singular arcs; an arbitrarily
small such translation preserves the strict mesh inequality. Tests
are chosen separately for each $R$, with the same numerical bounds;
no radius-independent exceptional null set is required.

The boundary of a backward continuity cell consists of cone-unstable
arcs and horizontal homogeneity boundaries, in lifted collision
coordinates. Such cells are convex along a cone-stable test: successive
crossings of their unstable boundaries separate different cells, and
an arc cannot end freely in a regular interior. This remains true
under inverse iteration, by invariance of the unstable cone and the
fact that new arc ends lie on previous singular or strip boundaries.
Consequently each connected test intersection has one homogeneous
inverse branch. Extra length subdivisions can only refine it; choose
one of them. Different $B$'s give disjoint inverse branches, so no
member is charged twice.

Finally $T_R^{m-1}A$ is obtained from $B$ by the last physical
one-step singularity cuts. There is no homogeneity cut at the final
image in the definition of $\mathcal M_0^{m,H}$. The one-step true
singularity set has a uniformly bounded number of lifted cone-stable
arcs, because only a fixed finite set of disks can be the next hit.
Each such arc crosses a cone-stable-convex backward cell in at most
one connected arc. Their ordered graph cuts make at most $D_*' +1$
components; take $D_*=D_*'+1$. This last bound counts physical
one-step cuts, not the countably many homogeneous strips.
\end{proof}

\begin{theorem}[Uniform bounded geometric pressure]
\label{thm:v100-uniform-pressure}
There is $Q_*<\infty$ such that
\begin{equation}\label{eq:v100-uniform-pressure}
             \sup_{R\in[R_-,R_+]}\sup_{m\ge0}Q_m(R)\le Q_*.
\end{equation}
In particular the actual geometric itinerary sums $P_m(R)$ in
\eqref{eq:v99-word-sum} satisfy the same bound.
\end{theorem}
\begin{proof}
There is a uniform one-step bound $J^sT_R\le C_1$. For every forward
cell $A$ counted in \eqref{eq:v100-long-cell}, put $B=B_{m-1}(A)$.
Dropping the last stable factor gives
\[
 \sup_AJ^sT_R^m\le
 C_1\sup_{y\in B\cap M_R'}
          J^sT_R^{m-1}(T_R^{-m+1}y).
\]
Assign $B$ to its first meeting test $W_j$. Compare a true stable
curve through $y$ to that test inside the same homogeneous backward
cell, using Lemma~\ref{lem:v100-geometric}. Inverting the two
backward Jacobians bounds the last supremum by
$C_d\sup_V J_VT_R^{m-1}$ for a representing homogeneous inverse
component $V$ of $W_j$. The distortion bound is uniform over the
whole cell; the test need not meet the particular final one-step
subcell $T_R^{m-1}A$.

There are at most $D_*$ such subcells per $B$. The injective
assignment in Lemma~\ref{lem:v100-tests} and the curve growth bound
therefore give, for $m\ge2$,
\[
 c_{\rm long}Q_m(R)
 \le C_1D_*C_d\sum_{j=1}^{N_*}
       \sum_{V\in\mathcal G_{m-1}(W_j)}\sup_VJ_VT_R^{m-1}
 \le C_1D_*C_dN_*C_g.
\]
The same geometric one-step estimates bound $Q_1$ uniformly, and
$Q_0=1$. Taking the larger of these constants proves the assertion.
Finally the homogeneous continuity partition refines each physical
itinerary; positivity gives $P_m(R)\le Q_m(R)$. Every sum is taken
at finite collision count before either supremum. In particular
\eqref{eq:v100-uniform-pressure} is not obtained by interchanging
$\sup_R$ with a fixed-table collision limit.
\end{proof}
